% Short running form; the full title is too wide for the head's right box.
\runninghead{The Garment the King Looks For}
\properlane{wedding-garment}{The Garment the King Looks For: the New Man Put On in Charity}
\label{sec:wedding-garment}

Heard from its end, the Gospel is about a man who is already inside. He was
called, he came, he sat down at the king's table, and he is found without the
one thing the king looks for. The Epistle of the same Mass has already named
that thing in its own language. ``Put on the new man, who according to God is
created in justice and holiness of truth'': the command is to clothe oneself,
and the Gospel's guest is condemned for not being clothed. Read together, the
two lessons are one demand seen from two sides. The Church is the wedding,
entry into it is by the king's call, and what the king looks for in those who
have come in is the new man put on, which is Christ, and the charity by which
Christ is put on. The Epistle shows that garment in four plain acts: truth to
the neighbour, anger ended before sunset, no place given to the devil, and work
that has something to give. And the Communion and Postcommunion pray for what
all this demands, that the one who has been commanded may be directed to keep
the commandments and made to cleave to them. St Gregory the Great and St
Augustine say what the garment is, and St Jerome says what it is made of; each
of them comments on the parable, and Jerome on the Epistle as well.

\subsection{Not baptism, not faith, not the altar: St Gregory and St Augustine}

St Gregory the Great preached on this Gospel to the people of Rome, and when he
reaches the man without a garment he begins by setting aside two answers that
lie close to hand. ``\latin{Si enim vestem nuptialem baptisma vel fidem dicimus,
quis sine baptismate et fide has nuptias intravit?}'' If we call the wedding
garment baptism or faith, who has entered this wedding without baptism and
faith? The unclothed guest is inside; whatever he lacks, it is not what brought
him in. ``\latin{Quid ergo debemus intelligere nuptialem vestem, nisi
charitatem?}'' What then should we understand by the wedding garment, if not
charity? And he gives a reason that turns from the guest to the bridegroom:
``\latin{Recte enim charitas nuptialis vestis vocatur, quia hanc in se conditor
noster habuit, dum ad sociandae sibi Ecclesiae nuptias venit}'': charity is
rightly called the wedding garment, because our Maker himself had it when he
came to the wedding at which he joined the Church to himself (\work{Homiliae in
Evangelia} 38, 9). The garment is the bridegroom's own dress. Whoever honours
him at his wedding wears what he wore to it.

Gregory then gives the garment a shape. It is woven, as cloth is woven, on two
beams, an upper and a lower, and so charity is held in two commandments, the
love of God and the love of neighbour (38, 10). It is tested by the hardest
case: \latin{Charitas autem vera est cum et in Deo diligitur amicus, et propter
Deum diligitur inimicus}, charity is true when the friend is loved in God and
the enemy is loved for God's sake (38, 11). Charity also explains the king's
greeting. The king calls the man \latin{Amice} and condemns him, and Gregory
reads the word as if the king had said more plainly: ``\latin{Amice, et non
amice; amice per fidem, sed non amice per operationem}'', friend and not
friend, friend by faith but no friend by what he has done (38, 12). The man's
silence follows from it. In that last rebuke, Gregory says, every argument of
excuse ceases, because the one who rebukes without is the witness of the
conscience that accuses within. To anyone who has the garment and has it
imperfectly he offers the psalm's hope, ``Thy eyes did see my imperfect being''
(Ps 138:16), and he turns his warning to the one who has none of it at all.

St Augustine makes the same elimination at greater length, in a sermon whose
heading calls it a sermon against the Donatists, on charity. The occasion
shapes the argument: the wedding hall holds good and bad together, and the
rite that brings them in is common to both, so that what separates them cannot
be a rite. ``Is it
Baptism? Without Baptism it is true no one attaineth to God; but not every one
that hath Baptism attaineth to Him. I cannot therefore understand Baptism, the
Sacrament itself that is, to be `the wedding garment;' for this garment I see
in the good, I see in the bad. Peradventure it is the Altar, or That which is
received at the Altar. But no; we see that many eat, and `eat and drink
judgment to themselves'.'' Not fasting, not coming to church, not miracles,
since the wicked do all these too (\work{Sermo} 90, 5). The garment is what the
good have and the bad have not, and Augustine finds it in St Paul: ``This is
the wedding garment: `Now the end of the commandment,' says the Apostle, `is
charity out of a pure heart, and of a good conscience, and of faith
unfeigned'\,'' (90, 6; 1 Tim 1:5). Not charity of any kind, he adds, since
thieves and the partisans of the chariot races love one another too; this
charity, from a pure heart.

Two further points in the sermon matter for the Mass. The first is where the
garment is worn. ``The garment that was looked for is in the heart, not on the
body; for had it been put on externally, it could not have been concealed even
from the servants'' (90, 4). The king sees what the servants who brought the
guests in did not see. The second is how it grows. ``Let charity be born in
thee, if it be yet unborn, and if it be born, be it nourished, fostered,
increased'' (90, 6). The garment is not received once and then possessed; it
is put on as charity advances, and Augustine ends the sermon by saying what
that advance restores: ``so be `the wedding garment' put on; so be the image
of God, after which we were created, by this our advancing, engraven anew in
us'' (90, 10). The Epistle of the Mass describes the new man as the one who is
created \latin{secundum Deum}, after God. Augustine does not cite the Epistle
there, but his last sentence is the Epistle's command seen from the inside:
the garment put on is the image of the Creator renewed.

\subsection{The garment of the new man: St Jerome and the Epistle}

St Jerome names the garment differently, and in doing so he joins the parable
to the Epistle's language himself. On the king's question he writes:
``\latin{Vestis autem nuptialis, praecepta sunt Domini, et opera quae
complentur ex lege et Evangelio, novique hominis efficiunt vestimentum}'', the
wedding garment is the Lord's precepts and the works fulfilled from the law and
the Gospel, which make the garment of the new man. The guest who is found on
the day of judgement \latin{sub nomine Christiano}, under the Christian name,
without the wedding garment, has not the garment of the heavenly man but a
soiled one, \latin{id est, veteris hominis exuvias}, the cast skin of the old
man (\work{Commentarii in Matthaeum} III, on Mt 22:11--12; PL 26,
cols.~160--161). The old man that the Epistle's readers were to put off
(Eph 4:22) is still on him.

Commenting on the Epistle itself, Jerome says who the new man is. He hears in
the command to put on the new man the same thing that Paul says elsewhere in
other words, \latin{Induite vos Christum Jesum} (Rom 13:14): ``\latin{Iste quippe est novus homo, quo universi credentes debemus
indui atque vestiri}'', this indeed is the new man, with whom all believers
ought to be clothed and dressed; and whoever imitates Christ's way of life, his
meekness and his humility, has put on the new man (\work{Commentarii in
Epistolam ad Ephesios} II, on 4:24; PL 26, col.~508). The two comments meet in
one image. The garment the king looks for is made of the Lord's precepts kept,
and the one who keeps them is clothed with Christ.

St Thomas Aquinas, commenting on the Epistle centuries later, reaches the
same identification by another road. As Adam was the first principle of
oldness in every thing, through whom sin entered into all, so, he writes,
``\latin{principium primum novitatis et renovationis Christus est}'', Christ is
the first principle of newness and renewal, and he at once adds Paul's
\latin{Induimini ergo Dominum nostrum Jesum Christum}. The Spirit of the mind in
which the Ephesians are to be renewed he takes first of the Holy Spirit
dwelling in the mind, the cause of renewal; the words \latin{qui creatus est}
he refers, in one of his readings, to Christ himself, formed in the Virgin's
womb by the Holy Spirit; and he ends his exposition of the verse with a
division that the rest of the lesson fills out, \latin{Vel ut sanctitas sit in
corde, veritas in ore, justitia in opere}, holiness in the heart, truth in the
mouth, justice in the deed (\work{In Epistolam ad Ephesios} c.~4, lect.~7).
Aquinas does the same with the parable. In his lectures on St Matthew he asks
what the garment is and answers, Christ. Some put Christ on in the sacrament,
some are in him by charity; and to have the wedding garment is to put on
Christ by good work, by a holy way of life and by true charity, so that if one
of these is wanting, it is evil (\work{Super Matthaeum} c.~22). The one Doctor thus
names Christ the new man of the Epistle and Christ the garment of the Gospel,
in two separate commentaries.

St John Chrysostom, preaching on the same verses of Ephesians, fixes what the
putting on changes: ``Seest thou that the subject is one, but the clothing is
twofold, that which is put off, and that which is put on?'' The man is the
same man; his clothing is exchanged, and the new man is created ``to be a son:
for this takes place from Baptism'' (\work{Homilies on Ephesians} 13, on
4:23--24). Schuster, reading this Mass's Epistle, says it in one sentence:
``This garment is Jesus Christ himself'' (p.~172). And the chapter from which
the Epistle is cut opens with a sentence that joins the call to the life it
asks for: ``I therefore, a prisoner in the Lord, beseech you that you walk
worthy of the vocation in which you are called'' (Eph 4:1). The call has been
given. The Epistle and the Gospel both ask what the called are wearing.

\subsection{What the new man does: truth, anger, and hands that give}

The Epistle does not leave the new man abstract. It names four acts, and the
commentators on the Epistle read them as the garment's visible side. St John
Chrysostom takes the first from the image Paul gives it, ``for we are members
one of another'': ``Let not the eye, saith he, lie to the foot, nor the foot to
the eye,'' since a body whose members deceive one another destroys itself.
On the second, ``Wouldest thou have thy fill of anger? One hour, or two, or
three, is enough for thee; let not the sun depart, and leave you both at
enmity. It was of God's goodness that he rose: let him not depart, having shone
on unworthy men.'' ``The blessed Paul dreads the night,'' he adds, lest anger
left burning in solitude be kindled again before morning. On the third: ``to be
at war with one another, is `to give place to the devil'\,''. And on the last he
sums up what the four have in common: ``Seest thou what are the members of the
old man? Falsehood, revenge, theft.'' The thief is not told merely to stop.
``If they labor, and charitably communicate to others, thus will they do away
the sin'' (\work{Homilies on Ephesians} 14). Aquinas reads the same verses with
the same point. The truth owed to the neighbour follows from membership,
\latin{Membra enim se invicem diligunt, et se juvant mutuo in veritate}, for
members love one another and help one another in truth; the first motion of
anger is allowed to human frailty, \latin{non permittitur mora}, its delay is
not; and labour is not only necessary but good, \latin{ut laborans possit
vivere, et Ut habeat unde tribuat necessitatem patienti}, that the labourer
may live and may have something to give (\work{In Epistolam ad Ephesios}
c.~4, lect.~8--9). Jerome is sharper still. ``\latin{Qui igitur ad hoc tantum
laborat, ut ipse non egeat, et a caeteris contrahit manum: quamvis applaudat
sibi, tamen Apostoli praeceptum non fecit}'': whoever works only so as not to
be in want himself, and draws back his hand from others, has not kept the
Apostle's precept, however he may applaud himself (on Eph 4:28, PL 26,
col.~512).

The hand that works in order to give is the Epistle's last image, and St
Gregory turns the Gospel's last image on the same hands. The king orders the
man bound hand and foot, and Gregory reads the binding as already begun in
this life: ``\latin{Pedes enim qui visitare aegrum negligunt, manus quae nihil
indigentibus tribuunt a bono opere jam ex voluntate ligatae sunt}'', feet that
neglect to visit the sick, and hands that give nothing to the needy, are
already bound from good work by their own will (\work{Homiliae in Evangelia}
38, 13). What the king does at the end, the guest has done to himself
beforehand. St Augustine, preaching on the parable in another sermon, gives the
same demand its positive form, and ties it back to the garment: ``Do not say;
`we are too poor to have that garment.' Clothe others, and ye are clothed
yourselves. It is winter, clothe the naked. Christ is naked; and He will give
you that `wedding garment' whosoever have it not. Run to Him, beseech Him; He
knoweth how to sanctify His faithful ones, He knoweth how to clothe His naked
ones.'' And then, with the binding in view: ``let not your works fail. If they
fail, with hands bound what canst thou do? with feet bound, whither wilt thou
fly?'' (\work{Sermo} 95, 7). The thief of the Epistle who works with his hands
so as to have something to give is weaving the wedding garment; the guest whose hands gave nothing arrives at the feast without it.

\subsection{Hands lifted at evening: St Hilary on the Gradual}

Between the Epistle and the Gospel the Mass sings, ``the lifting up of my
hands, as evening sacrifice''. St Hilary of Poitiers, expounding this verse,
distinguishes two positions of the hands. To spread the hands out is the
posture of prayer; to lift them up is something more. When Moses was told to
stand with his hands raised over the battle, the lifting showed \latin{excelsi
operis officium}, the office of a lofty work. And when St Paul bids men pray
``lifting up pure hands, without anger and contention'' (1 Tim 2:8), Hilary
observes, \latin{non habitum orandi constituit, sed praescriptum divinae
operationis adjecit}: he did not fix a posture of prayer but added a
prescription for divine work. Isaiah's God rejects those who spread out their
hands, not those who lift them, \latin{quia habitus orandi in manibus
extendendis est, perfecti autem operis efficientia est in elevandis}, because
the posture of praying lies in spreading the hands, but the doing of a
finished work lies in lifting them (\work{Tractatus super Psalmos} 140, 3).

Which evening sacrifice, then, pleases God? Hilary answers from the Gospel of
the last judgement: ``Come, ye blessed of my Father, possess you the kingdom
prepared for you \dots\ For I was hungry, and you gave me to eat: I was
thirsty, and you gave me to drink: I was a stranger, and you took me in''
(Mt 25:34--35). \latin{Manuum ergo elevatio, sacrificium vespertinum est. Non
enim sanguine et holocaustis nos, in quos consummatio saeculorum devenit,
sacrificamus Deo}: the lifting up of the hands, then, is the evening
sacrifice, for we upon whom the end of the ages has come do not sacrifice to
God with blood and holocausts (140, 4). The hands the Gradual lifts are hands
that feed, clothe, welcome and visit Christ in the least of his brethren.

St John Chrysostom, expounding the psalm's first verse in Greek before he
reaches its second, lays down the other condition Paul's sentence names. He
forbids prayer against enemies, after the example of the apostles and of
Stephen, and requires that the hands raised in prayer be raised without anger,
citing the same verse of the letter to Timothy (\work{Expositio in Psalmum}
140). St Robert Bellarmine cites it again on the Gradual's verse: ``the
Apostle, `Lifting up pure hands, without anger and strife'\,'' (on Ps 140:2).
Hilary's lifted hands are the Epistle's working hands, raised in prayer
without the anger the Epistle forbids to outlast the day; and they are the
opposite of the hands Gregory finds already bound, because they have given.

\subsection{What God commands is asked of God: the other elements}

The Communion states a command and answers it with a wish: ``Thou hast
commanded thy commandments to be kept most diligently. O! that my ways may be
directed to keep thy justifications.'' St Jerome said that the wedding garment
is \latin{praecepta Domini}, and the Communion's \latin{mandata} are those
precepts. Three expositors of the psalm read its wish as a prayer for grace.
St Augustine: ``When thou hearest, `O that,' recognise the words of one
wishing; and having recognised the expression of a wish, lay aside the pride of
presumption. \dots\ Therefore if man desireth what God chargeth, God must be
prayed to grant Himself what He enjoineth'' (\work{Enarrationes} 118, Aleph 5).
St Hilary, on the same verse, says that the justifications are many and that
\latin{nisi a Deo dirigamur, infirmes per naturam nostram erimus. Adjuvandi
igitur per gratiam ejus, dirigendique sumus}: unless we are directed by God we
shall be weak by our nature; we must therefore be helped and directed by his
grace (\work{Tractatus super Psalmos} 118, Aleph 12). And the Greek
\work{Expositiones in Psalmos} printed under the name of St Athanasius say of
the same verses that the psalmist, knowing that no one can keep the law without
help from above, secures its keeping for himself by prayer. The Postcommunion
then asks for exactly that help, that God's healing work \latin{tuis semper
faciat inhaerere mandatis}. Put together, the wedding garment is made of the
commandments kept, and the commandments are kept by the grace that the Mass,
at its end, asks for.

The other elements take their place around this centre. The Introit's Lord
bids his people attend \latin{legem meam}, and St Augustine, expounding the
psalm from which the verse comes, draws the distinction on which the question
of the garment turns: ``although all the Sacraments were common, grace, which is the
virtue of the Sacraments, was not common to all. \dots\ For even heretics have
the same Baptism, and false brethren too, in the communion of the Catholic
name'' (\work{Enarrationes} 77, 2). The guest without the garment is such a
brother. The Collect asks to be set free, in mind and body alike, to carry out
\latin{quae tua sunt}; the garment is put on by doing what is God's, and
Schuster calls the freedom the Collect asks for the liberty of the sons of
God. The Alleluia's praise is answered, in Augustine's reading of its verse,
by a God who hears those who love him: \latin{Exaudit quippe invocantem, quem
laudantem videt: laudantem videt, quem probat amantem}, he hears the one who
calls upon him whom he sees praising, and he sees praising the one whom he
proves to be loving (\work{Enarrationes} 104, 1). The one who is heard is the
one who loves, the guest in the garment of charity. The Offertory's walk
``in the midst of tribulation'' is the road on which the garment is kept. St
Augustine reads the verse as the condition of being revived at all:
``otherwise Thou wilt not revive me, unless I walk in the midst of
tribulation'' (\work{Enarrationes} 137, 10). And Chrysostom, in his homily on
the parable, says what the called must do once they are clothed: ``although to be called and to
be cleansed was of grace, yet, when called and clothed in clean garments, to
continue keeping them so, this is of the diligence of them that are called''
(\work{Homilies on Matthew} 69). The Secret asks that the gifts be
\latin{salutaria nobis}, and Schuster reads the prayer as a request for the
dispositions that let the oblation work its full effect in the soul
(p.~173); the garment is that disposition, brought to the altar.

\subsection{Gift or work: the strongest difficulty}

If entry into the wedding is the king's free call, and the guest is condemned
for lacking a garment, the call seems to demand what it did not give. The
witnesses answer the difficulty in different ways, and the differences are
real.

St Hilary makes the garment a gift. \latin{Vestitus autem nuptialis est gloria
Spiritus sancti, et candor habitus coelestis}: the wedding garment is the glory
of the Holy Spirit and the brightness of a heavenly dress, received in the
confession of a good answer at baptism and to be kept \latin{usque in coetum
regni coelorum immaculatus et integer}, spotless and whole until the assembly
of the kingdom of heaven (\work{Commentarius in Matthaeum} 22, 7; PL 9,
col.~1044). St Irenaeus, arguing against those who divided the God who calls
from the God who judges, says that ``we ought, after our calling, to be also
adorned with works of righteousness, so that the Spirit of God may rest upon
us; for this is the wedding garment'', and that those who were called to God's
supper but ``have not received the Holy Spirit, because of their wicked
conduct'' are cast out (\work{Against Heresies} IV.36.6). St Gregory and St
Augustine, as above, deny that the garment is baptism, because the unworthy
guest is baptized too. St John Chrysostom holds both sides in one sentence:
calling and cleansing are of grace; keeping the garment clean is the
diligence of the called.

The positions are narrower than they look. Hilary's garment is received in
baptism but must be kept spotless; Augustine's and Gregory's is charity, which a
baptized man may fail ever to put on. They agree that the guest is baptized and that
what he lacks is something a baptized man can lose or never wear, and they
differ over what to call it and where it comes from. Here the Communion's verse
gives the answer. The three expositors of that verse agree
that what God commands must be asked of him; and Hilary, who makes the garment
a gift of the Spirit on the Gospel, says on the psalm that the commandments
are kept only as God's grace directs. The garment is God's gift, and it must
be put on and kept; it is kept by praying for what is commanded. Augustine's
sermon puts the two sides into one exhortation: ``Christ is naked; and He will
give you that `wedding garment' \dots\ Run to Him, beseech Him.'' The
continuation of \work{The Liturgical Year} says the same in its own voice: the
king ``has given them abundant time to lay aside their tatters; and knowing,
that they could not get ready of themselves, he has placed at their disposal,
for the marriage-feast, the richest garments of his grace and virtues.''

A second difficulty is the one Hilary himself raises. How could a man swept in
from the highways be expected to have a garment? \latin{Numquid invitandorum
habitum designaverat? \dots\ quomodo unus omnibus poterat esse vestitus?} Had
the king prescribed a dress for those he invited? How could everyone have been
clothed alike? His answer is that the garment is inward, that human
simplicity hardly perceives a counterfeit mind, and that therefore God alone
found the unworthy guest (\work{Commentarius in Matthaeum} 22, 7). Augustine
answers in the same way, from the servants who did not notice: ``had it been
put on externally, it could not have been concealed even from the servants''.
The difficulty confirms that the garment is neither a piece of clothing nor a
rite.

\subsection{The four senses of this reading}

\begin{fourSenses}
\item[Literal.] Paul tells the Ephesians to be renewed in the spirit of their
mind and to put on the new man created according to God in justice and
holiness of truth, and he names what that means: truth spoken to the
neighbour, because they are members one of another; anger ended before the sun
sets; no place given to the devil; and work that can give to the needy. In the
parable a man who has come into the king's wedding without a wedding garment
is questioned, falls silent, and is bound and cast out; many are called, few
chosen.

\item[Allegorical.] The garment is Christ put on, the new and heavenly man
(Jerome; Aquinas, for whom Christ is the first principle of all newness as
Adam was of all oldness; Schuster), and the charity that the Bridegroom himself
wore when he came to join the Church to himself (Gregory). Its lack is the
cast skin of the old man (Jerome). The guest is a friend by faith and no friend
by his works (Gregory), and in putting on the garment the image of God after
which man was created is engraved anew (Augustine).

\item[Moral.] Charity toward God and neighbour, woven on both beams and proved
toward the enemy (Gregory; Augustine); truth among the members of one body,
anger settled the same day, and work done so as to give (the Epistle, with
Chrysostom, Aquinas and Jerome); hands lifted in the works of mercy and not
only spread in prayer (Hilary on the Gradual), since hands that give nothing
are already bound (Gregory); and what God commands, asked of God (Augustine and
Hilary on the Communion's verse).

\item[Anagogical.] The king's entry is the judgement, when the guest found under
the Christian name without the garment is rebuked (Jerome); what is sought at
the vineyard, at the building of the house and at the wedding feast is
\latin{non initia, sed finis}, not the beginning but the end (Jerome, on
Mt 22:14); no guest knows whether he is chosen, \latin{Quia vocati sumus,
novimus; si sumus electi, nescimus} (Gregory, \work{Homiliae} 38, 14); and the
garment received is to be kept spotless until the assembly of the kingdom of
heaven (Hilary).
\end{fourSenses}
